\documentclass[10pt,a4paper]{amsart}
\usepackage[margin=19mm]{geometry}
\usepackage{amsmath,amssymb,mathtools}
\usepackage[scr=rsfs,cal=euler]{mathalfa}
\usepackage{xfrac}
\usepackage[unicode,hidelinks,bookmarks=false]{hyperref}
\newcommand{\ie}{i.e.\ }
\newcommand{\cf}{cf.\ }
\newcommand{\resp}{respectively\ }
\newcommand{\Subsection}{\subsection}
\providecommand{\nobreakdash}{\nobreak\hbox{-}}
\setlength{\emergencystretch}{2em}
\multlinegap0pt
\usepackage{amssymb}
\usepackage[shortlabels]{enumitem}
\usepackage{xcolor}
\usepackage{tikz}
\usepackage{bm}
\newtheorem{claim}{Claim}
\newtheorem{proof*}{Proof*}
\usepackage{cleveref}
\crefname{claim}{Claim}{Claims}
\crefname{proof*}{Proof*}{Proof*s}
\title{Frustration index of a signed planar graph \\ and the feedback vertex set}
\author{Sirui Chen, Jiaao Li, Zhouningxin Wang}
\date{Source: arXiv:2607.17983v1 (CC BY 4.0)}
\begin{document}
\maketitle

\noindent\emph{Source and licence.} Frustration index of a signed planar graph \\ and the feedback vertex set. Version arXiv:2607.17983v1 (CC BY 4.0), distributed by arXiv under the Creative Commons Attribution 4.0 International licence, http://creativecommons.org/licenses/by/4.0/, as recorded in the arXiv licence metadata for this exact version and shown on the abstract page https://arxiv.org/abs/2607.17983v1. Full licence notice: \texttt{licenses/arxiv-2607.17983-CC-BY-4.0.txt}.\par\smallskip

\noindent\textit{Editorial scope.} Counted unit: the article's claim claim:order (source0.tex lines 516--518). The article's complete original proof of it is delivered with it (source0.tex lines 519--524, one \texttt{proof} environment of 6 lines). No mathematics is written, repaired or re-derived: the statement and the proof are the article's own lines, in the article's own notation, and no retained line is substituted or re-flowed.  No further source-local context is needed: every cross-reference of every retained line resolves inside the two delivered spans.  Nothing was omitted from any retained span, so the package carries no omission record.  No retained line contains a citation, so the wrapper carries no bibliography and no bibliography entry.  No retained line contains a figure environment, an image-inclusion command or drawing code, so no image is delivered and the package declares no asset dependency.  Editorial formatting only, in addition: an AMS \texttt{amsart} preamble that restates the article's own declarations and macros used by the retained lines, namely the environments \texttt{claim}, \texttt{proof*} on separate counters, together with the standard packages those lines need.  The retained environments print in this document as Claim 1, Proof* 1 in order of appearance, whereas the article prints the same items with the numbering of its own full document.  The complete native source of the article as distributed in the arXiv e-print for this exact version is bundled unchanged as \texttt{sources/205616/source0.tex}.
\begin{claim}\label{claim:order}
For each $i\in \{0,1,\ldots, s\}$, the graph $G_i-D_i$ admits a vertex ordering $\mathcal{O}_i=(v_1, v_2, \ldots, v_{n_i})$ such that for every $\ell\geq 2$, the number of edges joining $v_\ell$ to $\{v_1,\ldots,v_{\ell-1}\}$ is odd.  
\end{claim}

\begin{proof*}[Proof of \Cref{claim:order}]
We prove the statement by induction on $i$. The case $i=0$ is trivial. Assume that the statement holds for $i-1$, and let $\mathcal{O}_{i-1}$ be a vertex ordering of $G_{i-1}-D_{i-1}$ satisfying the required condition. In particular, the vertex $z$ to be split, either is the first vertex in the ordering, or has an odd number of incident edges to its earlier neighbors in $\mathcal{O}_{i-1}$.

We construct $\mathcal{O}_i$ from $\mathcal{O}_{i-1}$ as follows. If $z$ is the first vertex in $\mathcal{O}_{i-1}$, we replace $z$ by $z_1,z_2$. If $z_1$ is incident with an odd number of edges joining it to the earlier neighbors of $z$, then we replace $z$ in $\mathcal{O}_{i-1}$ by $z_1,z_2$. Otherwise, we replace $z$ by $z_2,z_1$.
We claim that $\mathcal{O}_i$ is a required vertex ordering of $G_i-D_i$. Indeed, recall that $z_1$ and $z_2$ are connected by an odd number of parallel edges in $G_i-D_i$. Exactly one of $z_1$ and $z_2$ is incident with an odd number of edges joining it to the earlier neighbors of $z$, and the other vertex is incident with an even number of such edges. We complete the proof of the claim.
\end{proof*}

\end{document}
